\section{Conceptual Framework and Justification for the Choice of Articles}
The M4 Competition article was selected because it provides validation of a hybrid model across a large panel of time series. While the first source offers an analysis of a single industry, the second source provides a true ``testing ground'' across 100{,}000 series, confirming the localized findings on a global scale.

\subsection{Validating the relevance of hybrid models}

The steel study develops a hybrid architecture, combining ARIMAX and Neural Networks, in order to integrate both time-series dynamics and exogenous factors into a single forecasting engine. One of the most significant findings of the M4 Competition was the victory of Slawek Smyl's ES-RNN model, described as a ``truly hybrid'' algorithm that combines Exponential Smoothing (ES) with Recurrent Neural Networks (RNN). The M4 Competition therefore demonstrates that the hybrid approach used in the steel study is not merely an industrial-specific phenomenon, but may instead represent a gold standard for future forecasting.

\subsection{Debunking the simple Machine Learning hype}

Both articles reach a similar conclusion: pure Machine Learning (ML) models often fail in time-series forecasting. The steel article notes that non-combined ML approaches lack the ability to capture all sources of uncertainty. Likewise, the M4 article reports the ``poor performances of the submitted pure ML methods,'' which were all less accurate than simple statistical benchmarks. This finding strongly supports the argument that combining statistical and machine learning approaches is more effective than relying on ML alone.

\subsection{Addressing uncertainty and risk}

Finally, both articles are relevant because of their shared focus on uncertainty management. The steel study emphasizes the need for hybrid models to account for model, parameter, and data uncertainties. The M4 Competition expands on this point by being the first competition to include \(95\%\) prediction intervals (PIs) in its evaluation, showing that hybrid models, such as Smyl's, are uniquely capable of specifying uncertainty with ``amazing precision.'' This is particularly important for the steel manufacturer's make-to-stock strategy, where understanding the range of possible demand is just as important as producing an accurate point forecast.

In summary, the second article acts as a global proof of concept for the specific industrial strategy proposed in the first, transforming a localized case study into a globally validated forecasting principle.